% Periodic content improvement; existing decay and nonvanishing are retained.
\section{Periodic content for the conductor-three frame}\label{sec:chi-periodic}

The translated-chain argument also applies to the $\chi$ frame. It improves
the content estimate without changing the node family, the decay certificate,
or the nonvanishing theorem of Section~\ref{sec:minor-nv}.

\begin{proposition}\label{prop:chi-periodic-local}
Let $p\ge5$ be prime, $p>2m$, and $p^2>3(r+2m)$. Suppose a cut-free
interval of the zone $[r,r+2m)$ contains $L$ consecutive nodes.
For every retained common-node set $J\subseteq[0,r)$, the coefficient
denominator of $P_J$ has $p$-exponent at most
\[
  2r+m-\max(0,L-m).
\]
Here the cut set consists of the singular indices and the points
$jp,r+jp$, $j\in\mathbf Z$.
\end{proposition}
\begin{proof}
The pole data are
\[
 \Delta_k=\frac1k+\frac6{3k-2}-\frac6{2k-1},\qquad
 D_k=\frac1{k^2}+\frac{18}{(3k-2)^2}-\frac{12}{(2k-1)^2},
 \qquad n=36.
\]
Thus $p\Delta_k,p^2D_k$ are integral, and both unscaled quantities are
integral away from $p\mid k(3k-2)(2k-1)$. This singular set is periodic
modulo $p$. After $x=U/p^2$, the scaled core is $UI+C$ with integral
constant entries. Its residue-class blocks have monic determinant, so
their cleared inverses have primitive denominators. Translation within
one cut-free cell identifies the corresponding core chains and their
cleared source vectors. The divided-difference formula then gives the
same antisymmetric first-order cancellation as in the uniform argument:
the cleared Schur entries inside the cell are integral, and entries
between distinct residues have denominator exponent at most one.
There are no repeated residues in the zone because $p>2m$.
The zero/one assignment bound is therefore
$2r+2m-L$ when $L>m$, and $2r+m$ otherwise.

The same bound holds after any integral constant perturbation $E$ of
the common core. Indeed, keep the old cleared partners $S_R,S_C$.
The residuals change by $-S_RE,-ES_C$, which are integral, and the
cleared second difference changes by $S_RES_C$, also integral.
The changed scaled core remains a monic shift of an integral matrix.
All nonnegative entry budgets, and hence the determinant bound, survive.
To recover a principal minor, add diagonal entries $0$ or $1$ only at
the omitted common coordinates and take successive determinant
differences. Multilinearity replaces those rows by identity rows;
deleting them leaves exactly the desired principal restriction.
This argument introduces no deletion penalty.
\end{proof}

For $p\equiv c\pmod6$, the singular endpoints, up to bounded integer
shifts, are $jp/6$ with
\[
 j\bmod6\in\{0,2,3\}\quad(c=1),\qquad
 j\bmod6\in\{0,3,4\}\quad(c=5).
\]
For consecutive endpoints $j,k$ from this dictionary, use the integral
witness interval
\[
 a=\max\!\left(r,\left\lfloor\frac{jp+4}{6}\right\rfloor+1\right),
 \qquad b=\min\!\left(r+2m,\left\lfloor\frac{kp}{6}\right\rfloor\right).
\]
The inequalities $jp+4<6t<kp$ exclude all three kinds of singularity;
the interval lies within one zone period. Also $2(b-a)\le p$ whenever
$b\ge a$. Thus a positive saving $(b-a-m)^+$ already implies $p>2m$.

The following disjoint windows give an affine lower bound
$Ar+Bp-100$ for the integer saving, uniformly for $r\le8m\le r+9$.
Only the level-one budget is $2r+m$ here; higher prime-power and bad-prime
costs are retained in the total denominator and contribute $o(r^2)$.
\begin{center}
\begin{tabular}{crrrrrr}
\toprule
$c$ & lower $p/r$ & upper $p/r$ & $A$ & $B$ & $j$ & $k$\\
\midrule
$1$ & $9/32$ & $2/7$ & $-9/8$ & $4$ & $21$ & $24$\\
$1$ & $2/7$ & $5/16$ & $-1/8$ & $1/2$ & $21$ & $24$\\
$1$ & $5/16$ & $9/28$ & $9/8$ & $-7/2$ & $21$ & $24$\\
$1$ & $3/8$ & $2/5$ & $-9/8$ & $3$ & $15$ & $18$\\
$1$ & $2/5$ & $5/12$ & $-1/8$ & $1/2$ & $15$ & $18$\\
$1$ & $5/12$ & $9/20$ & $9/8$ & $-5/2$ & $15$ & $18$\\
$1$ & $27/56$ & $1/2$ & $-9/8$ & $7/3$ & $12$ & $14$\\
$1$ & $1/2$ & $15/28$ & $-1/8$ & $1/3$ & $12$ & $14$\\
$1$ & $15/28$ & $9/16$ & $9/8$ & $-2$ & $12$ & $14$\\
$1$ & $9/16$ & $5/8$ & $-9/8$ & $2$ & $9$ & $12$\\
$1$ & $5/8$ & $2/3$ & $1/8$ & $0$ & $9$ & $12$\\
$1$ & $2/3$ & $3/4$ & $9/8$ & $-3/2$ & $9$ & $12$\\
$1$ & $3/4$ & $5/6$ & $-1/8$ & $1/6$ & $8$ & $9$\\
$1$ & $5/6$ & $27/32$ & $9/8$ & $-4/3$ & $8$ & $9$\\
\bottomrule
\end{tabular}
\end{center}
\begin{center}
\begin{tabular}{crrrrrr}
\toprule
$c$ & lower $p/r$ & upper $p/r$ & $A$ & $B$ & $j$ & $k$\\
\midrule
$5$ & $1/4$ & $5/18$ & $-1/8$ & $1/2$ & $24$ & $27$\\
$5$ & $5/18$ & $9/32$ & $9/8$ & $-4$ & $24$ & $27$\\
$5$ & $9/28$ & $1/3$ & $-9/8$ & $7/2$ & $18$ & $21$\\
$5$ & $1/3$ & $5/14$ & $-1/8$ & $1/2$ & $18$ & $21$\\
$5$ & $5/14$ & $3/8$ & $9/8$ & $-3$ & $18$ & $21$\\
$5$ & $3/8$ & $5/12$ & $-1/8$ & $1/3$ & $16$ & $18$\\
$5$ & $5/12$ & $27/64$ & $9/8$ & $-8/3$ & $16$ & $18$\\
$5$ & $9/20$ & $1/2$ & $-9/8$ & $5/2$ & $12$ & $15$\\
$5$ & $1/2$ & $9/16$ & $9/8$ & $-2$ & $12$ & $15$\\
$5$ & $9/16$ & $3/5$ & $-9/8$ & $2$ & $10$ & $12$\\
$5$ & $3/5$ & $5/8$ & $-1/8$ & $1/3$ & $10$ & $12$\\
$5$ & $5/8$ & $27/40$ & $9/8$ & $-5/3$ & $10$ & $12$\\
\bottomrule
\end{tabular}
\end{center}
Each row satisfies the explicit integer inequality
\[
 (b-a-m)^+\ge Ar+Bp-100.
\]
The prime number theorem in the two reduced classes modulo $6$ gives
the total gain
\[
 \frac12\sum_{\rm rows}\int_{\theta_-}^{\theta_+}(A+B\theta)\,d\theta
   =\frac{11983}{645120}.
\]
Subtracting this from the base content constant yields
\begin{equation}\label{eq:chi-periodic-content}
 \Lambda_\chi=\frac{10339}{3072}-\frac{11983}{645120}
   =\frac{2159207}{645120}=3.346985057\ldots .
\end{equation}
Proposition~\ref{prop:chi-periodic-local} makes this content estimate
uniform in $J$. The existing $D=3.7251$, $c=6.6510$ and prime interval
of width $2/5$ therefore give
\[
 \frac{c}{D-\Lambda_\chi}
   =\frac{536336640}{30491189}=17.589889328\ldots<17.60,
 \qquad D-\Lambda_\chi<\frac25.
\]
The periodic entry interface, cleared cancellation, integral perturbation
and principal-minor transfer, finite profile, prime summation and final
measure assembly are kernel checked in Lean.
